% 第7章 实证
\section{实证：尺度扫描与参照系分解}\label{sec:empirical}

本章把第 \ref{sec:measurement} 章与第 \ref{sec:reference} 章的命题在真实数据上跑一遍。三组实验各有分工：尺度扫描检验窗口与形状，参照系分解检验 $\alpha$ 路径，统计门槛的演示放在合成面板与真实数据的对照里。所有数字锚定同一份数据快照，脚本与本节的每个表一一对应。

\subsection{数据与装置}

按第 \ref{sec:criteria} 章流程的第一步，先把观测装置填清。

\textbf{数据底座}：本地通达信 vipdoc 日线（\texttt{.day}），窗口 2019-01-02 至 2026-10-08（$T=1881$）；universe 为沪深北 A 股 400 只（含北交所 15 只、窗口内上市 153 只），清单随快照冻结。价格乘后复权因子（第三方兜底来源，已对 400 只全量交叉核验，相对偏差中位数约 $2\times10^{-4}$），先于一切特征与收益计算。交易制度掩码按封板价精确判定涨跌停，新股上市后前 20 个交易日剔除；有效样本内封板合计 $1.96\%$（涨停 $1.49\%$、跌停 $0.48\%$）、一字板 $0.20\%$。快照校验和 \texttt{781cadd68dfbd170}。

\textbf{时间尺度}：日线分辨率，$\Delta_{\min}=1$ 日；信号在网格日收盘后计算、以收盘价成交（$\tau_e\approx0$ 的理想化口径）；观测噪声 $\nu_0$ 不单独估计，由检测门的统计量间接体现。

\textbf{特征与标签}：$K=8$ 的价量因子（1/5/20/60 日动量、20 日波动、量比、距 20 日均线偏离、Amihud 非流动性\cite{amihud2002}），逐期截面 $z$ 标准化，只用 $t$ 期及以前信息；标签为后复权收盘对收盘的下一日收益。

\textbf{成本}：三档分别记账——零成本、固定费率（佣金 $2.5$bp 与过户费 $0.1$bp 双边、卖出印花税按 2023-08-28 分段 $10$bp/$5$bp）、含冲击（逐股 ADV 取窗口内日成交额中位数，冲击率 $0.02\sqrt{\text{参与率}}$，参与率上限 $10\%$；口径与形式沿 Almgren 与 Chriss 2000\cite{almgrenchriss2000}、Almgren 等 2005\cite{almgren2005}）。单笔资金规模取 $10^8$ 元。

\subsection{实验一：尺度扫描}

操作定义与第 \ref{sec:measurement} 章一致：信号每日计算，但决策与调仓只发生在网格点 $\Delta\in\{1,2,3,5,10,20\}$（日），网格之间权重冻结；成本只在调仓日按真实换手收取。每个特征的方向取样本内更优者并在全部网格上固定（该选择计入试验账本，见第 \ref{sec:limits} 章的边界讨论）。

先看三个代表候选的完整曲线。它们分别对应三种形态：内点最优（量比）、右端存活（非流动性）、被成本整体压制（1 日反转）。

\begin{table}[htbp]
\centering
\caption{三个代表候选的尺度曲线（全样本；方向取样本内更优后固定）}
\label{tab:curves}
\begin{tabular}{@{}>{\raggedright\arraybackslash}p{2.6cm}r r r r@{}}
\toprule
候选（口径） & 网格 $\Delta$（日） & 年化换手 & 零成本夏普 & 含冲击夏普 \\
\midrule
量比（多空） & 1 & 154 & $3.422$ & $-0.355$ \\
              & 2 & 96 & $2.712$ & $0.074$ \\
              & 3 & 71 & $2.262$ & $0.176$ \\
              & 5 & 47 & $2.007$ & $\mathbf{0.483}$ \\
              & 10 & 27 & $1.420$ & $0.412$ \\
              & 20 & 15 & $0.763$ & $0.124$ \\
\midrule
非流动性（多空） & 1 & 14 & $1.221$ & $1.014$ \\
              & 2 & 10 & $1.206$ & $1.030$ \\
              & 3 & 9 & $1.194$ & $1.031$ \\
              & 5 & 8 & $1.166$ & $1.017$ \\
              & 10 & 6 & $1.164$ & $1.027$ \\
              & 20 & 5 & $1.245$ & $\mathbf{1.131}$ \\
\midrule
1 日反转（多空） & 1 & 322 & $1.294$ & $-7.690$ \\
              & 2 & 160 & $1.328$ & $-3.238$ \\
              & 3 & 107 & $1.111$ & $-2.131$ \\
              & 5 & 65 & $0.466$ & $-1.502$ \\
              & 10 & 32 & $0.198$ & $-0.896$ \\
              & 20 & 16 & $-0.153$ & $-0.698$ \\
\bottomrule
\end{tabular}
\end{table}

三条曲线的形状差异，正是第 \ref{sec:measurement} 章两组命题要刻画的东西。量比的多空曲线是标准的内点最优：左端 $1$ 日尺度上，年化换手 $154$ 倍，含冲击夏普被成本压到 $-0.355$；放宽到 $5$ 日，换手降到 $47$ 倍，含冲击夏普冲上 $0.483$；再放宽到 $20$ 日，零成本夏普已从 $3.422$ 衰减到 $0.763$（信息衰减开始主导），含冲击夏普回落到 $0.124$。成本摊薄与信息衰减两股力，一个在左端抬着走、一个在右端拉着走，中间夹出峰值。非流动性是另一个极端：零成本夏普几乎平坦（$1.16$ 到 $1.25$ 之间，信号衰减极慢），含冲击夏普随网格放宽单调改善，在 $20$ 日处达到最高 $1.131$——对慢信号，窗口的右端一直开着。1 日反转则是反面教材：零成本夏普在 $1$--$2$ 日最高（$1.29$、$1.33$，信息新鲜时最值钱），随后快速衰减，到 $20$ 日已转负；含冲击夏普在所有网格上为负——\textbf{它的窗口是被成本整体压死的：信息没错、衰减也有，但每一项都被 $322$ 倍的年换手吃光。}

把全部 16 种构建的含冲击夏普排成矩阵，可以看到形状与信号快慢的对应关系。

\begin{table}[htbp]
\centering
\caption{全部构建的含冲击夏普（列：网格 $\Delta$，单位日；峰位指含冲击夏普最大的网格）}
\label{tab:matrix}
\begin{tabular}{@{}>{\raggedright\arraybackslash}p{3.1cm}r r r r r r >{\raggedright\arraybackslash}p{2.0cm}@{}}
\toprule
构建 & 1 & 2 & 3 & 5 & 10 & 20 & 峰位 \\
\midrule
1 日反转（多空） & $-7.69$ & $-3.24$ & $-2.13$ & $-1.50$ & $-0.90$ & $-0.70$ & 右端，全负 \\
1 日反转（仅多头） & $-2.98$ & $-1.22$ & $-0.60$ & $-0.40$ & $-0.03$ & $0.15$ & 右端 \\
5 日动量（多空） & $-0.88$ & $-0.64$ & $-0.41$ & $-0.60$ & $-0.17$ & $-0.14$ & 右端，全负 \\
5 日动量（仅多头） & $-0.65$ & $-0.46$ & $-0.17$ & $-0.16$ & $0.04$ & $0.21$ & 右端 \\
20 日动量（多空） & $0.46$ & $0.52$ & $\mathbf{0.55}$ & $0.30$ & $0.10$ & $0.09$ & 3 \\
20 日动量（仅多头） & $0.15$ & $0.23$ & $0.32$ & $0.28$ & $0.32$ & $0.43$ & 右端 \\
60 日动量（多空） & $0.80$ & $\mathbf{0.86}$ & $0.70$ & $0.66$ & $0.51$ & $0.33$ & 2 \\
60 日动量（仅多头） & $0.46$ & $0.50$ & $\mathbf{0.50}$ & $0.49$ & $0.47$ & $0.49$ & 3 \\
波动率（多空） & $\mathbf{0.73}$ & $0.72$ & $0.61$ & $0.62$ & $0.43$ & $0.24$ & 1 \\
波动率（仅多头） & $0.67$ & $0.67$ & $0.64$ & $\mathbf{0.67}$ & $0.62$ & $0.60$ & 5（近平） \\
量比（多空） & $-0.36$ & $0.07$ & $0.18$ & $\mathbf{0.48}$ & $0.41$ & $0.12$ & 5 \\
量比（仅多头） & $-0.17$ & $0.06$ & $0.19$ & $0.32$ & $0.38$ & $0.43$ & 右端 \\
距均线偏离（多空） & $0.50$ & $\mathbf{0.54}$ & $0.51$ & $0.38$ & $0.08$ & $-0.07$ & 2 \\
距均线偏离（仅多头） & $0.11$ & $0.15$ & $\mathbf{0.29}$ & $0.24$ & $0.19$ & $0.29$ & 3 \\
非流动性（多空） & $1.01$ & $1.03$ & $1.03$ & $1.02$ & $1.03$ & $\mathbf{1.13}$ & 右端 \\
非流动性（仅多头） & $0.89$ & $0.91$ & $0.90$ & $0.91$ & $0.92$ & $\mathbf{1.00}$ & 右端 \\
\bottomrule
\end{tabular}
\end{table}

表 \ref{tab:matrix} 的读法有四层。

\textbf{第一，内点最优是常态而非例外。}多空口径下有 4 个构建出现清晰的内点峰（20 日动量在 $3$ 日、60 日动量在 $2$ 日、量比在 $5$ 日、距均线偏离在 $2$ 日），仅多头口径另有 2 个（60 日动量在 $3$ 日、距均线偏离在 $3$ 日）。峰的位置集中在 $2$--$5$ 日——这正好落在“成本已经摊薄、信息尚未衰尽”的地带。

\textbf{第二，两端的情形同样有信息。}最慢的信号（非流动性，零成本夏普几乎不随网格变化）在右端吃满成本摊薄的好处，峰值在测试范围的最右；最快的信号（1 日反转、5 日动量）来不及被任何网格救活，多空口径全矩阵为负。快慢之分在表里转化为窗口位置的差别，与第 \ref{sec:measurement} 章“窗口位置由衰减速率的量级决定”方向一致。

\textbf{第三，尺度把“可交易集合”整体扩容。}在 $\Delta=1$ 的单一网格上（下一小节的视角），16 种构建里只有 2 种含冲击为正；把网格放宽后，正的构建增加——20 日动量、60 日动量、波动率、量比的多头化口径都在中段网格转正。\textbf{“哪些候选可交易”这个问句的答案，与观测尺度绑定；这就是观测缝隙理论最直接的经验含义。}

\textbf{第四，不要读过头。}这些是全样本口径的描述性曲线：未做逐（候选，网格）的滚动样本外检验（已有的样本外核验在单特征与随机候选上做过，见下一小节）；方向选择含样本内成分；数据只有日频，$\Delta<1$ 的一侧不可见。曲线的价值在于形状与理论相容、峰位有秩序，而不是某个具体网格上的具体数值。（“右端”指峰值出现在测试网格的最大值 $20$ 日处，真实峰位可能更右，受测试范围所限。）

\subsection{实验二：单网格视角与它的局限}

上一小节末尾提到，$\Delta=1$ 的单一网格会给出一个更“悲观”的可交易集合。把这个视角单独展开有意义，因为它对应大量实证研究的默认做法。在 $1$ 日网格上，16 种构建里含冲击夏普为正的只有 2 种，都来自非流动性（多空 $1.014$、仅多头 $0.890$）；其余候选的换手跨度从 $25$ 倍到 $322$ 倍不等，量比多空 $154$ 倍换手、含冲击 $-0.355$，1 日反转多空 $322$ 倍、含冲击 $-7.690$。若据此下结论“成本后几乎无人生还”，结论本身没错——错的是把 $1$ 日网格的结论当成与尺度无关的结论。第 \ref{sec:measurement} 章的命题在这里显示它的用处：\textbf{“成本后能否存活”不是候选的固有属性，而是候选与观测网格的联合属性；换一个网格，存活的集合就变。}

\subsection{实验三：参照系分解}

尺度扫描说的是“以什么频率看”，这一小节换一个维度：把候选的毛收益在参照工具逐步扩张的链条上做嵌套回归，看 $\alpha$ 的路径。链条为：无参照（仅现金）$\to$ 加等权全市场 $\to$ 加非流动性多空工具 $\to$ 加 1 日反转多空工具。四件工具自身在三步中的表现先记下：等权全市场对现金的年化 $\alpha$ 为 $+0.121$；非流动性工具对现金 $+0.118$、加等权后 $+0.114$（对市场近乎中性）；1 日反转工具对现金 $-0.118$、加等权后 $-0.106$——作为一条可交易组合，它在样本期内的绝对收益为负，这是它的客观属性，不影响它作为参照工具的功能。

\begin{table}[htbp]
\centering
\caption{参照系分解链：候选的年化 $\alpha$（全样本、零成本毛收益；方向取样本内更优）}
\label{tab:decomp}
\begin{tabular}{@{}>{\raggedright\arraybackslash}p{3.4cm}r r r r@{}}
\toprule
候选（方向） & 对现金 & $+$等权市场 & $+$非流动性工具 & $+$1 日反转工具 \\
\midrule
非流动性（仅多头，$+$） & $0.287$ & $0.166$ & $-0.013$ & $-0.005$ \\
量比（仅多头，$-$） & $0.189$ & $0.066$ & $0.063$ & $0.050$ \\
1 日反转（仅多头，$-$） & $0.081$ & $-0.051$ & $-0.049$ & $-0.148$ \\
\texttt{rand10}（仅多头，$+$） & $0.175$ & $0.055$ & $0.066$ & $0.028$ \\
\bottomrule
\end{tabular}
\end{table}

表 \ref{tab:decomp} 里能读出三层信息。

\textbf{第一，市场暴露是最大的单一吞并源。}加一件等权市场工具，四个候选的年化 $\alpha$ 从 $(0.287,0.189,0.081,0.175)$ 变成 $(0.166,0.066,-0.051,0.055)$——多数腰斩以上，1 日反转候选甚至转负。这正是实务里“先剥 beta 再看 alpha”的标准动作在分解链上的位置：它对应公式 \eqref{eq:decomp} 里 $\gamma$ 接近 $1$ 而 $\alpha_{\mathcal{F}}(g)$ 为市场的历史溢价那一项。

\textbf{第二，吞并是选择性的，且量与公式的乘积吻合。}非流动性候选在被同源工具吞并时给出了最干净的演示：$\alpha$ 从 $0.166$ 掉到 $-0.013$，下降 $0.179$；而它对非流动性工具的暴露为 $1.564$，工具自身的 $\alpha$ 为 $0.114$，乘积 $1.564\times0.114\approx0.178$——与下降量在三位小数内吻合。同一件工具对 \texttt{rand10} 与量比几乎没有影响（$0.055\to0.066$、$0.066\to0.063$），因为它们对非流动性工具的暴露都接近于零（$-0.096$、$0.030$）。参照系只吞“与它同源”的 $\alpha$，这一步是公式的数值注脚。

\textbf{第三，第三件工具上的两次乘积同样成立。}加 1 日反转工具后：\texttt{rand10} 的 $\alpha$ 掉 $0.038$，对应暴露 $-0.337$ 乘工具 $\alpha$ $-0.113$，乘积 $0.038$；1 日反转候选掉 $0.100$，对应暴露 $-0.884$ 乘 $-0.113$，乘积 $0.100$。符号与量都对得上——\textbf{分解公式不是记账修辞，它是可以被逐个乘数检验的等式。}

落到结论上，最重要的一行是非流动性：它的正 $\alpha$ 与“非流动性因子自身的溢价”在数据上不可分——加入同源工具后归零。第 \ref{sec:limits} 章会把这条线索的其余检验（子集衰减、规模中性化、逐年稳定性）摊开，得到的结论与此处一致：**它的收益更像一种风险定价，而不是一条等待提取的缝隙。**参照系分解的价值，在这里从“改数字”升级为“定性”——它把一条线索从候选名单里请了出去，理由可以被任何持有该工具数据的人复算。

\subsection{统计门槛的对照演示}

第 \ref{sec:criteria} 章把统计校正列为独立一道门。它的必要性可以用两组现成结果对照说明，一组来自合成面板，一组来自本章的真实数据。

合成面板（$T=1200$、$N=60$、纯噪声主导）上，按样本内夏普从 $61$ 个候选里挑选，选中的是一条纯噪声线性组合：样本内夏普 $1.196$，滚动前向样本外夏普 $0.477$，缩水夏普（标准误口径）$0.053$，回测过拟合概率 $0.567$——统计量集体亮红灯，这个候选被第四道门拦下。真实数据上，同一套流程选中的随机组合（\texttt{rand10}）样本内 $3.649$、样本外 $2.671$、缩水夏普 $0.9991$、过拟合概率 $0.000$——它通过了统计门；但它的年化换手 $182$ 倍，含冲击夏普 $-1.492$（成本占毛收益 $1.49$ 倍），倒在第二道门。至于单特征量比的 $t$ 值（$-9.03$）在多重检验下仍然过硬，且对序列相关稳健（Newey--West 调整后几乎不动；非流动性由 $3.91$ 降到 $3.45$，仍过门槛\cite{newyewest1987}）——它倒在第一、二道门的组合上（高换手，且多空口径依赖做空）。

这组对照要说明的是四道门的\textbf{独立性}：每个候选卡在不同位置，没有哪一道可以被另一道代替。统计上最结实的候选（\texttt{rand10}）被成本拦下；成本上最结实的候选（非流动性）在参照系与可提取性上存疑（第 \ref{sec:limits} 章展开）。\textbf{“过了哪几道门”比“是不是好策略”更接近可操作的语言。}

\subsection{口径自检与复现}

两处自检要交代。其一，$\Delta=1$ 应与既有独立审计的逐日调仓口径对齐：多空构建的换手与两档夏普精确复现（非流动性：$14$ 倍、$1.221$、$1.014$；1 日反转：$322$ 倍与原表符号翻转后的绝对值一致）；仅多头构建存在 $\le0.003$ 的差异，来源是组合归一化与不可成交约束的先后次序，两个口径都合理，不影响任何结论方向。其二，本节的实验数据曾在快照冻结后被本地数据库更新（新增了更晚的日线），为对齐原始快照，脚本按固定截止日截断并读取冻结的 400 只清单复算，结果与快照时代逐字节一致。

复现信息如下：
\begin{itemize}
  \item 复现命令：\texttt{python code/\allowbreak gap\_scan.py}、\texttt{python code/\allowbreak ref\_decomp.py}；
  \item 数据层：\texttt{code/\allowbreak real\_data\_frozen.py}；
  \item 输出文件：\texttt{outputs/\allowbreak gap\_scan.txt}、\texttt{outputs/\allowbreak ref\_decomp.txt}；
  \item 快照校验和：\texttt{781cadd68dfbd170}；复权因子缓存与 universe 清单见 \texttt{outputs/} 目录。
\end{itemize}

\subsection{本章要点}

尺度扫描给出了窗口与形状的直接证据：内点最优普遍存在于中速信号，慢信号的窗口右端全开，快信号与成本赛跑失败；单网格视角会系统性低估可交易集合。参照系分解把“$\alpha$ 相对什么”变成一条可读的路径。两组对照共同支撑本章的主句：\textbf{可行性由装置与成本共同划定，而装置的两个维度——尺度与参照系——都可以在数据上被逐项检验，不必停留在口头声明。}